% !TEX root = ../../sheet.tex


%%%%%%%%%%%%%%%%%%%%%%%%%%%%%%%%%%%%%%%%%%%%%%%%%%%%%%%%%%%%%%%%%%%%%%%%%%%%%%%%%%%%%%%%%%%%%%%%%%%%
\section{Entwurfsmuster} %%%%%%%%%%%%%%%%%%%%%%%%%%%%%%%%%%%%%%%%%%%%%%%%%%%%%%%%%%%%%%%%%%%%%%%%%%%
%%%%%%%%%%%%%%%%%%%%%%%%%%%%%%%%%%%%%%%%%%%%%%%%%%%%%%%%%%%%%%%%%%%%%%%%%%%%%%%%%%%%%%%%%%%%%%%%%%%%


\blindtext[1]

\begin{tasks}
% Task (a)
\item Schreiben Sie eine Funktion \fsharp|replicate: 'a -> Nat -> 'a list|, die einen beliebigen Wert \fsharp|x| sowie eine
	natürliche Zahl \fsharp|n| nimmt und eine Liste zurückgibt, die \fsharp|n| mal das Element \fsharp|x| enthält.\\[5pt]
	Beispiele:\\[5pt]
	\fsharp|replicate 0N 3N = [0N; 0N; 0N]|\\
	\fsharp|replicate 3N 0N = []|\\
	\fsharp|replicate "Hallo" 2N = ["Hallo"; "Hallo"]|
\begin{tasks}
	% Subtask (i)
	\item Gehen Sie \textbf{strikt nach dem Peano Entwurfsmuster} vor.
		\ifthenelse{\boolean{lsg}}{%
			\originalinput{%
				rangestartafterstring = //replicatePeanoBegin,
				rangestopbeforestring = //replicatePeanoEnd,
			}{D:/Bachelorarbeit/24-ba-jakob-meyerhoefer/evaluation/template/Implementation.fs}
		}{%
			\originalinput{%
				rangestartafterstring = //replicatePeanoBegin,
				rangestopbeforestring = //replicatePeanoEnd,
			}{D:/Bachelorarbeit/24-ba-jakob-meyerhoefer/evaluation/template/Implementation.fs}
		}
		\solutioninput{%
			rangestartafterstring = //replicatePeanoBegin,
			rangestopbeforestring = //replicatePeanoEnd,
		}{D:/Bachelorarbeit/24-ba-jakob-meyerhoefer/evaluation/template/Solution.fs}
	% Subtask (ii)
	\item Gehen Sie \textbf{strikt nach dem Leibniz Entwurfsmuster} vor.
		\ifthenelse{\boolean{lsg}}{%
			\originalinput{%
				rangestartafterstring = //replicateLeibnizBegin,
				rangestopbeforestring = //replicateLeibnizEnd,
			}{D:/Bachelorarbeit/24-ba-jakob-meyerhoefer/evaluation/template/Implementation.fs}
		}{%
			\originalinput{%
				rangestartafterstring = //replicateLeibnizBegin,
				rangestopbeforestring = //replicateLeibnizEnd,
			}{D:/Bachelorarbeit/24-ba-jakob-meyerhoefer/evaluation/template/Implementation.fs}
		}
		\solutioninput{%
			rangestartafterstring = //replicateLeibnizBegin,
			rangestopbeforestring = //replicateLeibnizEnd,
		}{D:/Bachelorarbeit/24-ba-jakob-meyerhoefer/evaluation/template/Solution.fs}
\end{tasks}
% Task (b)
\pagebreak\item Schreiben Sie eine Funktion \fsharp|exists|, die eine Funktion \fsharp|f| vom Typ \fsharp|Nat -> bool| sowie zwei natürliche Zahlen
	\fsharp|lower| und \fsharp|upper| nimmt und prüft, ob es eine Zahl \fsharp|x| im Intervall \fsharp|lower..upper| gibt, für die \fsharp|f x = true| ist.\\[5pt]
	Beispiele:\vspace{-.5\baselineskip}
	\begin{multicols}{2}
		\noindent			
		\fsharp|exists (fun x -> x = 9N) 10N 20N = false|\\
		\fsharp|exists (fun x -> x = 10N) 10N 20N = true|\\
		\fsharp|exists (fun x -> x = 15N) 10N 20N = true|\\
		\fsharp|exists (fun x -> x = 20N) 10N 20N = true|\\
		\fsharp|exists (fun x -> x = 21N) 10N 20N = false|\\
		\fsharp|exists (fun x -> true) 20N 10N = false|\\
		\fsharp|exists (fun x -> x = 42N) 42N 42N = true|
	\end{multicols}
	\ifthenelse{\boolean{lsg}}{%
			\originalinput{%
				rangestartafterstring = //existsBegin,
				rangestopbeforestring = //existsEnd,
			}{D:/Bachelorarbeit/24-ba-jakob-meyerhoefer/evaluation/template/Implementation.fs}
		}{%
			\originalinput{%
				rangestartafterstring = //existsBegin,
				rangestopbeforestring = //existsEnd,
			}{D:/Bachelorarbeit/24-ba-jakob-meyerhoefer/evaluation/template/Implementation.fs}
		}
		\solutioninput{%
			rangestartafterstring = //existsBegin,
			rangestopbeforestring = //existsEnd,
		}{D:/Bachelorarbeit/24-ba-jakob-meyerhoefer/evaluation/template/Solution.fs}
% Task (c)
\item Schreiben Sie eine rekursive Funktion \fsharp|mod7: Nat -> Nat|, die für eine gegebene Zahl \fsharp|n| den Rest der Ganzzahldivision durch \fsharp|7| (also \fsharp|n| modulo \fsharp|7|) berechnet.
Gehen Sie strikt nach dem Peano-Entwurfsmuster vor, verwenden Sie nicht die Operationen \fsharp|%| oder \fsharp|/|.\\[5pt]
	Beispiele:\vspace{-.5\baselineskip}
	\begin{multicols}{3}
		\noindent
		\fsharp|mod7 0N = 0N|\\
		\fsharp|mod7 1N = 1N|\\
		\fsharp|mod7 2N = 2N|\\
		\fsharp|mod7 3N = 3N|\\
		\fsharp|mod7 4N = 4N|\\
		\fsharp|mod7 5N = 5N|\\
		\fsharp|mod7 6N = 6N|\\
		\fsharp|mod7 7N = 0N|\\
		\fsharp|mod7 8N = 1N|
	\end{multicols}
	\ifthenelse{\boolean{lsg}}{%
		\originalinput{%
			rangestartafterstring = //mod7Begin,
			rangestopbeforestring = //mod7End,
			highlightlines = {2,5},
			highlightcolor = red!30,
		}{D:/Bachelorarbeit/24-ba-jakob-meyerhoefer/evaluation/template/Implementation.fs}
	}{%
		\originalinput{%
			rangestartafterstring = //mod7Begin,
			rangestopbeforestring = //mod7End,
		}{D:/Bachelorarbeit/24-ba-jakob-meyerhoefer/evaluation/template/Implementation.fs}
	}
	\solutioninput{%
		rangestartafterstring = //mod7Begin,
		rangestopbeforestring = //mod7End,
		highlightlines = {2,5},
		highlightcolor = green!30,
	}{D:/Bachelorarbeit/24-ba-jakob-meyerhoefer/evaluation/template/Solution.fs}
% Task (d)
\pagebreak\itemtodo TODO\\[5pt]
	Beispiele:\vspace{-.5\baselineskip}
	\begin{multicols}{5}
		\noindent
		\fsharp|mod5 0N = 0N|\\
		\fsharp|mod5 1N = 1N|\\
		\fsharp|mod5 2N = 2N|\\
		\fsharp|mod5 3N = 3N|\\
		\fsharp|mod5 4N = 4N|\\
		\fsharp|mod5 5N = 0N|\\
		\fsharp|mod5 19N = 4N|\\
		\fsharp|mod5 20N = 0N|\\
		\fsharp|mod5 21N = 1N|\\
		\fsharp|mod5 22N = 2N|
	\end{multicols}
	\begin{solutionwithcomment}{%
		\blindtext[1]
		\begin{itemize}[nosep]
			\item A
			\item B
			\item C
		\end{itemize}
	}
	let f a b =
		let result = a + b
		result
	\end{solutionwithcomment}
\end{tasks}


%%%%%%%%%%%%%%%%%%%%%%%%%%%%%%%%%%%%%%%%%%%%%%%%%%%%%%%%%%%%%%%%%%%%%%%%%%%%%%%%%%%%%%%%%%%%%%%%%%%%
% EOF %%%%%%%%%%%%%%%%%%%%%%%%%%%%%%%%%%%%%%%%%%%%%%%%%%%%%%%%%%%%%%%%%%%%%%%%%%%%%%%%%%%%%%%%%%%%%%
%%%%%%%%%%%%%%%%%%%%%%%%%%%%%%%%%%%%%%%%%%%%%%%%%%%%%%%%%%%%%%%%%%%%%%%%%%%%%%%%%%%%%%%%%%%%%%%%%%%%